\subsubsection{Pattern checking}\label{sec:pattern-checking}

\begin{figure}
\setlength{\jot}{2pt}
\begin{subfigure}{\textwidth}
\flushleft \shadebox{$\Sigma \vdash C[\seq{\sigma}] \instanceAt \tau$}
\begin{mathpar}
\small
\inferrule*[lab={\ruleName{instance-above}}]
{
  \Sigma \vdash \tau \subty C[\seq{\sigma}] \\
  \seq{\sigma}' \tyEquiv \seq{\sigma} \text{ whenever } \Sigma \vdash \tau \subty C[\seq{\sigma}']
}
{
  \Sigma \vdash C[\seq{\sigma}] \instanceAt \tau
}
\and
\inferrule*[lab={\ruleName{instance-below}}]
{
  \Sigma \vdash C[\seq{\sigma}] \subty \tau \\
  \seq{\sigma}' \tyEquiv \seq{\sigma} \text{ whenever } \Sigma \vdash C[\seq{\sigma}'] \subty \tau
}
{
  \Sigma \vdash C[\seq{\sigma}] \instanceAt \tau
}
\end{mathpar}
\subcaption{$C[\seq{\sigma}]$ is the instance of $C$ at $\tau$}
\label{fig:instance}
\end{subfigure}

\vspace{2mm}

\begin{subfigure}{\textwidth}
\flushleft \shadebox{$\Sigma; \Gamma \vdash p \checks \tau \binding \Delta$}
\begin{mathpar}
\small
\inferrule*[lab={\ruleName{check-var}}]
{
  \Gamma(x) = \Unbound
}
{
  \Sigma; \Gamma \vdash x \checks \tau \binding \set{\bind{x}{\tau}}
}
\and
\inferrule*[lab={\ruleName{check-wild}}]
{
  \strut
}
{
  \Sigma; \Gamma \vdash \patWild \checks \tau \binding \varnothing
}
\and
\inferrule*[lab={\ruleName{check-literal}}]
{
  \strut
}
{
  \Sigma; \Gamma \vdash \ell \checks \tau \binding \varnothing
}
\and
\inferrule*[lab={\ruleName{check-as}}]
{
  \Gamma(x) = \Unbound \\
  \Sigma; \Gamma \vdash p \checks \tau \binding \Delta
}
{
  \Sigma; \Gamma \vdash p\;\pyinline{as}\;x \checks \tau \binding \Delta \disjunion \set{\bind{x}{\matchedAs{\Sigma}(\Gamma, \tau, p)}}
}
\and
\inferrule*[lab={\ruleName{check-union}}]
{
  \Sigma; \Gamma \vdash p \checks \sigma \binding \Delta \\
  \Sigma; \Gamma \vdash p \checks \tau \binding \Delta'
}
{
  \Sigma; \Gamma \vdash p \checks \tyUnion{\sigma}{\tau} \binding \set{\bind{x}{\joinTy{\Sigma}{\Delta(x)}{\Delta'(x)}} \mid x \in \binds(p)}
}
\and
\inferrule*[lab={\ruleName{check-tuple}}]
{
  \Sigma; \Gamma \vdash p_i \checks \tau_i \binding \Delta_i \quad (\forall i)
}
{
  \Sigma; \Gamma \vdash \patTuple{p_1, \ldots, p_n} \checks \tyTuple{\tau_1, \ldots, \tau_n} \binding \Delta_1 \disjunion \cdots \disjunion \Delta_n
}
\and
\inferrule*[lab={\ruleName{check-tuple-none}}]
{
  \tau \text{ has no list}\added{\text{ or range}}\text{ values} \\
  \text{no disjunct of } \tau \text{ a tuple type of length } |\seq{p}|
}
{
  \Sigma; \Gamma \vdash \patTuple{\seq{p}} \checks \tau \binding \set{\bind{x}{\tyNever} \mid x \in \binds(\patTuple{\seq{p}})}
}
\and
\inferrule*[lab={\ruleName{check-list}}]
{
  \Sigma; \Gamma \vdash p_i \checks \tau \binding \Delta_i \quad (\forall i)
}
{
  \Sigma; \Gamma \vdash \patList{\seq{p}} \checks \tyList{\tau} \binding \Delta_1 \disjunion \cdots \disjunion \Delta_n
}
\and
\inferrule*[lab={\ruleName{check-list-none}}]
{
  \tau \text{ has no tuple}\added{\text{ or range}}\text{ values} \\
  \text{no disjunct of } \tau \text{ a list type}
}
{
  \Sigma; \Gamma \vdash \patList{\seq{p}} \checks \tau \binding \set{\bind{x}{\tyNever} \mid x \in \binds(\patList{\seq{p}})}
}
\and
\inferrule*[lab={\ruleName{check-dict}}]
{
  \Sigma; \Gamma \vdash p_i \checks \tau \binding \Delta_i \quad (\forall i)
}
{
  \Sigma; \Gamma \vdash \patDict{\bind{\seq{w}}{\seq{p}}} \checks \tyDict{\tau} \binding \Delta_1 \disjunion \cdots \disjunion \Delta_n
}
\and
\inferrule*[lab={\ruleName{check-dict-none}}]
{
  \text{no disjunct of } \tau \text{ a dictionary type}
}
{
  \Sigma; \Gamma \vdash \patDict{\bind{\seq{w}}{\seq{p}}} \checks \tau \binding \set{\bind{x}{\tyNever} \mid x \in \binds(\patDict{\bind{\seq{w}}{\seq{p}}})}
}
\and
\inferrule*[lab={\ruleName{check-constr}}]
{
  \Gamma \vdash q : \Type{C} \\
  \Sigma \vdash C[\seq{\sigma}] \instanceAt \tau \\
  \tyInst{\fields{\Sigma}(C)}{\seq{\sigma}} = \bind{\seq{x}}{\seq{\sigma}'} \\
  \fieldMap{\Sigma}(C, a) = \set{\bind{\seq{x}}{\seq{p}}} \\
  \Sigma; \Gamma \vdash p_i \checks \sigma'_i \binding \Delta_i \quad (\forall i)
}
{
  \Sigma; \Gamma \vdash \patConstr{q}{a} \checks \tau \binding \Delta_1 \disjunion \cdots \disjunion \Delta_n
}
\and
\inferrule*[lab={\ruleName{check-constr-none}}]
{
  \Gamma \vdash q : \Type{C} \\
  \Sigma \nvdash \tau \subty C[\seq{\sigma}] \text{ and } \Sigma \nvdash C[\seq{\sigma}] \subty \tau \quad (\forall \seq{\sigma})
}
{
  \Sigma; \Gamma \vdash \patConstr{q}{a} \checks \tau \binding \set{\bind{x}{\tyNever} \mid x \in \binds(\patConstr{q}{a})}
}
\end{mathpar}
\subcaption{Pattern $p$ checks against type $\tau$, binding $\Delta$}
\label{fig:checking-pat}
\end{subfigure}
\caption{Pattern checking}
\label{fig:pattern-checking}
\end{figure}


A pattern must \emph{check} against the remaining type, $\Sigma; \Gamma \vdash p \checks \tau \binding \Delta$ (\figref{checking-pat}), before it is matched against the residual (\ruleName{match-case}). The bindings $\Delta$ give the types of the pattern variables. The judgement holds when every tuple pattern in $p$ is at a type which has no list values, every list pattern is at a type which has no tuple values, and every constructor pattern is at a type with at most one comparable instantiation of its class, up to equivalence. List values and tuple values are defined by subtyping:
\begin{defalign}
\tau \text{ has list values} & {}\iff \Sigma \vdash \tyList{\sigma} \subty \tau \text{ for some } \sigma \text{, or a disjunct of } \tau \text{ is a type variable} \\
\tau \text{ has tuple values} & {}\iff \Sigma \vdash \tyTuple{\seq{\sigma}} \subty \tau \text{ for some } \seq{\sigma} \text{, or a disjunct of } \tau \text{ is a type variable}
\end{defalign}
So list types, $\tySized$ and $\tyObject$ have list values; tuple types, $\tySized$ and $\tyObject$ have tuple values; and a type variable, which may stand for any type, has both. The rules check a sub-pattern at the type of its position, a tuple component, list element, dictionary value or field, and a pattern at a union at each of its disjuncts. A variable pattern binds at the type at which it is checked, and a pattern named by \pyinline{as} is bound at $\matchedAs{\Sigma}(\Gamma, \tau, p)$; the bindings of sub-patterns compose by disjoint union, so a pattern that binds a name twice is ill-formed, and at a union the bindings of the disjuncts are joined. A pattern also checks against a type which has no disjunct of the form matched by the pattern: a tuple pattern when no disjunct is a tuple type of the pattern's length, a list pattern when no disjunct is a list type, a dictionary pattern when no disjunct is a dictionary type, or a constructor pattern when its class has no comparable instantiation (\ruleName{check-tuple-none}, \ruleName{check-list-none}, \ruleName{check-dict-none}, \ruleName{check-constr-none}). Such a pattern matches no shape of the type and binds its variables at $\tyNever$, the unit of the join. At one disjunct of a union it therefore leaves the values of that disjunct to later cases, and a case with a pattern matching no shape at all is unreachable.

The conditions on tuple and list patterns ensure that pattern matching always has a result. Under the evaluation rules, matching a tuple pattern against a list or a list pattern against a tuple is undefined (\S\ref{aux:match}), although Python permits both. Both forms of pattern therefore fail at $\tyList{\tyInt} \mathbin{|} \tyTuple{\tyInt, \tyInt}$, at $\tySized \mathbin{|} \tyTuple{\tyInt, \tyInt}$, and at a type variable.

A constructor pattern cannot name $\clsRange$, because Python does not support positional sub-patterns
on that class (\ruleName{pat-constr}). The condition on constructor patterns supplies the type arguments omitted by the pattern: they are recovered by the judgement $\Sigma \vdash C[\seq{\sigma}] \instanceAt \tau$ from the type $\tau$ at which the pattern is checked, the \emph{instance} of $C$ at $\tau$ being the instantiation of $C$ above $\tau$, or the instantiation below $\tau$, when there is one up to equivalence (\ruleName{instance-above}, \ruleName{instance-below}). Below $C$ the base classes of the type determine the instance: at \pyinline{IntLeaf}, a subclass of \pyinline{Leaf[int]}, the instance of \pyinline{Leaf} is \pyinline{Leaf[int]}. Above $C$ the base arguments of $C$ determine it: at \pyinline{Tree[int]} the instance of \pyinline{Leaf} is \pyinline{Leaf[int]}, because \pyinline{Leaf[T]} extends \pyinline{Tree[T]}, whereas at \tyObject{} a class with type parameters has no instance, since every instantiation lies below \tyObject.
The instance judgement is decided as follows. A pattern is checked at the disjuncts of a union separately
(\ruleName{check-union}), so the type $\tau$ at which an instance is taken is never a union. If $\tau$ is an
instantiated class $D[\seq{\sigma}]$ with $C$ an ancestor of $D$, the instance is the instantiation of $C$ reached
from $D[\seq{\sigma}]$ along the base classes, which is unique because type arguments are invariant
(\ruleName{instance-above}). If $\tau$ is $D[\seq{\sigma}]$ with $D$ an ancestor of $C$, the instantiations of $C$
below $\tau$ are the $C[\seq{\tau}]$ for which the instantiation of $D$ reached from $C[\seq{\tau}]$ along the base
classes has arguments equivalent to $\seq{\sigma}$, and Lemma~\ref{lem:instantiation-candidates} bounds the instantiations
to be examined to finitely many instantiation candidates: one instantiation candidate up to equivalence gives the instance
(\ruleName{instance-below}), none means that $C$ has no instance at $\tau$, and two that are not equivalent mean
that $C$ has an instantiation comparable with $\tau$ but no instance, so that a constructor pattern at $\tau$ is
ill-formed. At \tyObject{} and at \tyNever{} every instantiation of $C$ is comparable with $\tau$, so $C$ has an
instance only if it has no type parameters. At any other type $C$ has no instantiation comparable with $\tau$, and no instance.
\begin{specdefinition}{Instantiation candidate}
Suppose $D \in \ancestors{\Sigma}(C)$. An \emph{instantiation candidate} for $C$ at $D[\seq{\sigma}]$ is an instantiation $C[\seq{\tau}]$ in
which each $\tau_i$ is the join of a subset of $\subterms(\seq{\sigma}) \cup \set{\tyObject}$
(\S\ref{aux:subterms}).
\end{specdefinition}

\begin{lemma}[Sufficiency of instantiation candidates]
\label{lem:instantiation-candidates}
Suppose $D \in \ancestors{\Sigma}(C)$. If any instantiation of $C$ is below
$D[\seq{\sigma}]$, then some instantiation candidate is below $D[\seq{\sigma}]$; and if two instantiations of $C$ that are not equivalent
are below $D[\seq{\sigma}]$, then two instantiation candidates below $D[\seq{\sigma}]$ are not equivalent.
\end{lemma}

\subsubsection{Shapes and shape typing}\label{sec:shapes}

\begin{figure}
\small
\setlength{\jot}{2pt}
\begin{subfigure}{\textwidth}
\begin{tabular}[t]{rcll}
\grammarGroup{Heads} \\
$h$ & $::=$ & $\ell$ & (Literal) \\
    & $\mid$ & $C$ & (Class) \\
    & $\mid$ & $\headList{n}$ & (List of length $n$) \\
$H$ & $::=$ & $\set{h_1, \ldots, h_n}$ & (Excluded heads)
\end{tabular}
\hfill
\begin{tabular}[t]{rcll}
\grammarGroup{Shapes} \\
$k$ & $::=$ & $\shapeAny{\tau}{H}$ & (Rest) \\
    & $\mid$ & $\shapeAny{\patConstr{C[\seq{\tau}]}{k_1, \ldots, k_n}}{H}$ & (Constr) \\
    & $\mid$ & $\patTuple{k_1, \ldots, k_n}$ & (Tuple) \\
    & $\mid$ & $\shapeList{\tau}{k_1, \ldots, k_n}$ & (List) \\
    & $\mid$ & $\shapeDict{\tau}{\beta}{H}$ & (Dict) \\
$K$, $L$ & $::=$ & $\set{k_1, \ldots, k_n}$ & (Shape set)
\end{tabular}

\subcaption{Heads and shapes}
\label{fig:shapes-syntax}
\end{subfigure}

\vspace{2mm}

\begin{subfigure}{\textwidth}
\flushleft \shadebox{$\Sigma \vdash h : \tau$}
\begin{mathpar}
\small
\inferrule*[lab={\ruleName{head-literal}}]
{
  \Sigma \vdash \tyLiteral{\ell} \subty \tau
}
{
  \Sigma \vdash \ell : \tau
}
\and
\inferrule*[lab={\ruleName{head-literal-float}}]
{
  n \text{ a float literal} \\
  \Sigma \vdash \tyFloat \subty \tau
}
{
  \Sigma \vdash n : \tau
}
\and
\inferrule*[lab={\ruleName{head-class}}]
{
  \Sigma \vdash C[\seq{\sigma}] \subty \tau
}
{
  \Sigma \vdash C : \tau
}
\and
\inferrule*[lab={\ruleName{head-list}}]
{
  \strut
}
{
  \Sigma \vdash \headList{n} : \tyList{\tau}
}
\end{mathpar}
\subcaption{Head $h$ has type $\tau$}
\label{fig:head-typing}
\end{subfigure}

\vspace{2mm}

\begin{subfigure}{\textwidth}
\flushleft \shadebox{$\Sigma \vdash k : \tau$}
\begin{mathpar}
\small
\inferrule*[lab={\ruleName{shape-rest}}]
{
  \Sigma \vdash h : \tau \quad (\forall h \in H) \\
  \tau \text{ neither union nor dictionary type}
}
{
  \Sigma \vdash \shapeAny{\tau}{H} : \tau
}
\and
\inferrule*[lab={\ruleName{shape-constr}}]
{
  \tyInst{\fields{\Sigma}(C)}{\seq{\tau}} = \bind{\seq{x}}{\seq{\sigma}} \\
  \Sigma \vdash h : C[\seq{\tau}] \quad (\forall h \in H) \\
  \Sigma \vdash k_i : \sigma_i \quad (\forall i)
}
{
  \Sigma \vdash \shapeAny{\patConstr{C[\seq{\tau}]}{\seq{k}}}{H} : C[\seq{\tau}]
}
\and
\inferrule*[lab={\ruleName{shape-tuple}}]
{
  \Sigma \vdash k_i : \tau_i \quad (\forall i)
}
{
  \Sigma \vdash \patTuple{\seq{k}} : \tyTuple{\seq{\tau}}
}
\and
\inferrule*[lab={\ruleName{shape-list}}]
{
  \Sigma \vdash k_i : \tau \quad (\forall i)
}
{
  \Sigma \vdash \shapeList{\tau}{\seq{k}} : \tyList{\tau}
}
\and
\inferrule*[lab={\ruleName{shape-sub}}]
{
  \Sigma \vdash k : \sigma \\
  \Sigma \vdash \sigma \subty \tau
}
{
  \Sigma \vdash k : \tau
}
\and
\inferrule*[lab={\ruleName{shape-dict}}]
{
  \Sigma \vdash h : \tyStr \quad (\forall h \in H) \\
  \Sigma \vdash \beta(w) : \tau \quad (\forall w \in \dom(\beta)) \\
  \dom(\beta) \cap H = \varnothing
}
{
  \Sigma \vdash \shapeDict{\tau}{\beta}{H} : \tyDict{\tau}
}
\end{mathpar}
\subcaption{Shape $k$ has type $\tau$}
\label{fig:shape-typing}
\end{subfigure}
\caption{Shapes and shape typing}
\label{fig:shapes}
\end{figure}


A shape $k$ (\figref{shapes})
denotes a set of values of its type $\shapeTy(k)$ (\S\ref{aux:shape-type}). The form $\shapeAny{\tau}{H}$ denotes the
values of type $\tau$ whose outermost head is not among those in $H$; the other forms are constructor,
tuple, list and dictionary shapes, each with shapes in place of sub-patterns. A list shape also records its
element type, which the element shapes do not determine. A constructor shape also carries an excluded
set: $\shapeAny{\patConstr{C[\seq{\tau}]}{\seq{k}}}{H}$ denotes the instances of $C[\seq{\tau}]$, or of any subclass not below a
class of $H$, whose fields, those of $C$, have shapes $\seq{k}$. A rest shape at a class type denotes
the same values as the constructor shape with a rest shape at each field, and is expanded into that
form only by splitting (\figref{split}): a class may refer to itself through its fields, so expanding
in advance would not terminate.

A head itself denotes a set of values: a literal its value, a class its instances, and $\headList{n}$
the lists of length $n$. Head typing $\Sigma \vdash h : \tau$ holds when the values of $h$ lie within
$\tau$. Shape typing $\Sigma \vdash k : \tau$ holds when $\shapeTy(k)$ is a subtype of $\tau$ and $k$ is
well-formed: each excluded head is typed at the type of the shape (at $\tyStr$, for the keys of a
dictionary shape) and each component shape at the type of that component.

The auxiliary function $\shapes{\Sigma}(\tau, H)$ (\S\ref{aux:shapes}) gives the shapes of $\tau$
once the heads in $H$ are excluded. A union contributes the shapes of its disjuncts and a dictionary type
the dictionary form; every other type contributes the single shape $\shapeAny{\tau}{H}$, or no shape at
all when $H$ exhausts $\tau$. The rest form therefore never carries a union or dictionary type, and
\ruleName{shape-rest} states this restriction as a side condition.

\subsubsection{Matching}\label{sec:matching}

\begin{figure}
\small
\setlength{\jot}{2pt}
\flushleft \shadebox{$\Sigma; \Gamma \vdash k, p \matches K, L$}
\begin{mathpar}
\small
\inferrule*[lab={\ruleName{pat-var}}]
{
  \strut
}
{
  \Sigma; \Gamma \vdash k, x \matches \set{k}, \varnothing
}
\and
\inferrule*[lab={\ruleName{pat-wild}}]
{
  \strut
}
{
  \Sigma; \Gamma \vdash k, \patWild \matches \set{k}, \varnothing
}
\and
\inferrule*[lab={\ruleName{pat-as}}]
{
  \Sigma; \Gamma \vdash k, p \matches K, L
}
{
  \Sigma; \Gamma \vdash k, p\;\pyinline{as}\;x \matches K, L
}
\and
\inferrule*[lab={\ruleName{pat-split}}]
{
  \Sigma; \Gamma \vdash k, p \splits K, L \\
  \Sigma; \Gamma \vdash K, p \matches K', L'
}
{
  \Sigma; \Gamma \vdash k, p \matches K', L \disjunion L'
}
\and
\inferrule*[lab={\ruleName{pat-literal}}]
{
  \strut
}
{
  \Sigma; \Gamma \vdash \shapeAny{\tyLiteral{\ell}}{\varnothing}, \ell \matches \set{\shapeAny{\tyLiteral{\ell}}{\varnothing}}, \varnothing
}
\and
\inferrule*[lab={\ruleName{pat-literal-float}}]
{
  \added{n \text{ a float literal}} \\
  \added{n \notin H} \\
  \added{\Sigma \vdash \tyFloat \subty \tau}
}
{
  \added{\Sigma; \Gamma \vdash \shapeAny{\tau}{H}, n \matches \set{\shapeAny{\tau}{H}}, \set{\shapeAny{\tau}{H \cup \set{n}}}}
}
\and
\inferrule*[lab={\ruleName{pat-tuple}}]
{
  \Sigma; \Gamma \vdash \seq{k}, \seq{p} \matches K, L
}
{
  \Sigma; \Gamma \vdash \patTuple{\seq{k}}, \patTuple{\seq{p}} \matches \set{\patTuple{\seq{k}\,'} \mid \seq{k}\,' \in K}, \set{\patTuple{\seq{l}} \mid \seq{l} \in L}
}
\and
\inferrule*[lab={\ruleName{pat-list}}]
{
  \Sigma; \Gamma \vdash \seq{k}, \seq{p} \matches K, L
}
{
  \Sigma; \Gamma \vdash \shapeList{\tau}{\seq{k}}, \patList{\seq{p}} \matches \set{\shapeList{\tau}{\seq{k}\,'} \mid \seq{k}\,' \in K}, \set{\shapeList{\tau}{\seq{l}} \mid \seq{l} \in L}
}
\and
\inferrule*[lab={\ruleName{pat-dict}}]
{
  \seq{w} \text{ distinct} \\
  k_i = \beta(w_i) \quad (\forall i) \\
  \Sigma; \Gamma \vdash \seq{k}, \seq{p} \matches K, L
}
{
  \Sigma; \Gamma \vdash \shapeDict{\tau}{\beta}{H}, \patDict{\bind{\seq{w}}{\seq{p}}} \matches \set{\shapeDict{\tau}{\beta \override \set{\bind{\seq{w}}{\seq{k}\,'}}}{H} \mid \seq{k}\,' \in K}, \set{\shapeDict{\tau}{\beta \override \set{\bind{\seq{w}}{\seq{l}}}}{H} \mid \seq{l} \in L}
}
\and
\inferrule*[lab={\ruleName{pat-constr}}]
{
  \Gamma \vdash q : \Type{D} \\
  D \neq \clsRange \\
  D \in \ancestors{\Sigma}(C) \\
  \fields{\Sigma}(D) = \scheme{\seq{\alpha}}{\bind{\seq{x}}{\seq{\sigma}}} \\
  \fieldMap{\Sigma}(D, a) = \set{\bind{\seq{x}}{\seq{p}}} \\
  p'_i = \patWild \quad (\forall i) \\
  \Sigma; \Gamma \vdash \seq{k}, \seq{p} \mathbin{\concat} \seq{p}\,' \matches K, L
}
{
  \Sigma; \Gamma \vdash \shapeAny{\patConstr{C[\seq{\tau}]}{\seq{k}}}{H}, \patConstr{q}{a} \matches \set{\shapeAny{\patConstr{C[\seq{\tau}]}{\seq{k}\,'}}{H} \mid \seq{k}\,' \in K}, \set{\shapeAny{\patConstr{C[\seq{\tau}]}{\seq{l}}}{H} \mid \seq{l} \in L}
}
\end{mathpar}

\vspace{2mm}

\flushleft \shadebox{$\Sigma; \Gamma \vdash \seq{k}, \seq{p} \matches K, L$}
\begin{mathpar}
\small
\inferrule*[lab={\ruleName{pat-nil}}]
{
  \strut
}
{
  \Sigma; \Gamma \vdash \varepsilon, \varepsilon \matches \set{\varepsilon}, \varnothing
}
\and
\inferrule*[lab={\ruleName{pat-cons}}]
{
  \Sigma; \Gamma \vdash k, p \matches K, L \\
  \Sigma; \Gamma \vdash \seq{k}, \seq{p} \matches K', L'
}
{
  \Sigma; \Gamma \vdash k \cons \seq{k}, p \cons \seq{p} \matches K \times K', (L \times \set{\seq{k}}) \disjunion (K \times L')
}
\end{mathpar}

\vspace{2mm}

\flushleft \shadebox{$\Sigma; \Gamma \vdash L, p \matches K, L'$}
\begin{mathpar}
\small
\inferrule*[lab={\ruleName{pat-shapes}}]
{
  L' = \set{k \in L \mid \Sigma; \Gamma \vdash k, p \matches K_k, L_k} \\
  L' \neq \varnothing
}
{
  \Sigma; \Gamma \vdash L, p \matches \textstyle\bigdisjunion_{k \in L'} K_k, \; \textstyle\bigdisjunion_{k \in L'} L_k \disjunion (L \setminus L')
}
\end{mathpar}
\caption{Pattern $p$ matches shape $k$, giving matched shapes $K$ and residual $L$; lifted to sequences and ordered sets of shapes}
\label{fig:typing-pat}
\end{figure}


\emph{Matching} (\figref{typing-pat}) has three forms. The form on a single shape $k$ matches a pattern against that
shape, giving matched shapes $K$ and residual $L$. The form
on sequences matches the sub-patterns of a tuple, list, dictionary or constructor pattern against the
component shapes position by position. The form on
an ordered set of shapes $L$, used by the match rules, matches each shape of $L$ in turn. The pattern need not match every shape of $L$, so \ruleName{pat-shapes} takes the shapes of $L$
matched by the pattern, of which there must be at least one, and passes the unmatched shapes into the
residual untouched.

\begin{figure}
\small
\setlength{\jot}{2pt}
\flushleft \shadebox{$\Sigma; \Gamma \vdash k, p \splits K, L$}
\begin{mathpar}
\small
\inferrule*[lab={\ruleName{split-literal}}]
{
  \ell \notin H \\
  \Sigma \vdash \tyLiteral{\ell} \subty \tau \\
  \tau \text{ not a literal type}
}
{
  \Sigma; \Gamma \vdash \shapeAny{\tau}{H}, \ell \splits \set{\shapeAny{\tyLiteral{\ell}}{\varnothing}}, \shapes{\Sigma}(\tau, H \cup \set{\ell})
}
\and
\inferrule*[lab={\ruleName{split-tuple}}]
{
  |\seq{\tau}| = |\seq{p}|
}
{
  \Sigma; \Gamma \vdash \shapeAny{\tyTuple{\seq{\tau}}}{\varnothing}, \patTuple{\seq{p}} \splits \set{\patTuple{\seq{k}} \mid \seq{k} \in \shapes{\Sigma}(\seq{\tau})}, \varnothing
}
\and
\inferrule*[lab={\ruleName{split-list}}]
{
  n = |\seq{p}| \\
  \headList{n} \notin H
}
{
  \Sigma; \Gamma \vdash \shapeAny{\tyList{\tau}}{H}, \patList{\seq{p}} \splits \set{\shapeList{\tau}{\seq{k}} \mid \seq{k} \in \shapes{\Sigma}(\tau, \varnothing)^n}, \shapes{\Sigma}(\tyList{\tau}, H \cup \set{\headList{n}})
}
\and
\inferrule*[lab={\ruleName{split-dict}}]
{
  w_i \notin \dom(\beta) \cup H \\
  w_j \in \dom(\beta) \quad (\forall j < i)
}
{
  \Sigma; \Gamma \vdash \shapeDict{\tau}{\beta}{H}, \patDict{\bind{\seq{w}}{\seq{p}}} \splits \set{\shapeDict{\tau}{\beta \override \set{\bind{w_i}{k}}}{H} \mid k \in \shapes{\Sigma}(\tau, \varnothing)}, \set{\shapeDict{\tau}{\beta}{(H \cup \set{w_i})}}
}
\and
\inferrule*[lab={\ruleName{split-class}}]
{
  \Gamma \vdash q : \Type{C} \\
  \ancestors{\Sigma}(C) \cap H = \varnothing \\
  \Sigma \vdash C[\seq{\tau}'] \instanceAt \tau \\
  \meetTy{\Sigma}{\tau}{C[\seq{\tau}']} = D[\seq{\sigma}] \\
  \tyInst{\fields{\Sigma}(D)}{\seq{\sigma}} = \bind{\seq{x}}{\seq{\sigma}'} \\
  H' = \set{h \in H \mid \Sigma \vdash h : D[\seq{\sigma}]}
}
{
  \Sigma; \Gamma \vdash \shapeAny{\tau}{H}, \patConstr{q}{a} \splits \set{\shapeAny{\patConstr{D[\seq{\sigma}]}{\seq{k}}}{H'} \mid \seq{k} \in \shapes{\Sigma}(\seq{\sigma}')}, \shapes{\Sigma}(\tau, H \cup \set{C})
}

\and
\inferrule*[lab={\ruleName{split-subclass}}]
{
  \Gamma \vdash q : \Type{D} \\
  \Sigma \vdash D[\seq{\sigma}] \instanceAt C[\seq{\tau}] \\
  \Sigma \vdash D[\seq{\sigma}] \subty C[\seq{\tau}] \\
  D \neq C \\
  \ancestors{\Sigma}(D) \cap H = \varnothing \\
  \tyInst{\fields{\Sigma}(D)}{\seq{\sigma}} = \tyInst{\fields{\Sigma}(C)}{\seq{\tau}} \mathbin{\concat} \bind{\seq{y}}{\seq{\sigma}'} \\
  H' = \set{h \in H \mid \Sigma \vdash h : D[\seq{\sigma}]}
}
{
  \Sigma; \Gamma \vdash \shapeAny{\patConstr{C[\seq{\tau}]}{\seq{k}}}{H}, \patConstr{q}{a} \splits \set{\shapeAny{\patConstr{D[\seq{\sigma}]}{\seq{k} \mathbin{\concat} \seq{k}\,'}}{H'} \mid \seq{k}\,' \in \shapes{\Sigma}(\seq{\sigma}')}, \set{\shapeAny{\patConstr{C[\seq{\tau}]}{\seq{k}}}{(H \cup \set{D})}}
}
\end{mathpar}
\caption{Shape $k$ splits at the head tested by $p$, giving shapes $K$ and residual $L$}
\label{fig:split}
\end{figure}


A pattern matches a shape directly if the shape has the head tested by the pattern. Otherwise the shape
is first \emph{split} by the auxiliary judgement $\Sigma; \Gamma \vdash k, p \splits K, L$ (\figref{split}):
the shapes $K$ have the head tested by $p$, with the shapes of the component types drawn from $\shapes{\Sigma}$,
and the residual $L$ is the shape $k$ with that head excluded as well. The \ruleName{pat-split} rule matches
the pattern against $K$ and adds $L$ to the residual. Only the forms
with an excluded set split, and no head types at a tuple type, so the excluded set of a tuple shape is
empty and \ruleName{split-tuple} leaves the empty residual. Literal types do not include floats, so a float
literal pattern matches a rest shape directly (\ruleName{pat-literal-float}) rather than splitting it; the matched shape
is the whole shape, and the residual is the shape with the literal excluded. The matched shape therefore includes values
which are not equal to the literal. Those values are also in the residual, so the residual of an enclosing pattern is the
same as if the matched shape contained only the literal.
A constructor pattern whose class lies below a
class of the excluded set matches nothing, because an earlier case for the superclass has already
matched every instance. A constructor pattern matches the instance of its class at the type of the shape (\figref{instance}): a pattern \pyinline{Leaf(v)} at \pyinline{Tree[int]} matches \pyinline{Leaf[int]} and binds \pyinline{v} at \pyinline{int}.

A dictionary shape has a separate form,
$\shapeDict{\tau}{\beta}{H}$, standing for the dictionaries that hold each key $w$ in the domain of
$\beta$ with a value of shape $\beta(w)$, but that hold no key of $H$. A dictionary pattern takes its keys in order, with no key occurring twice. A key in
neither $\beta$ nor $H$ splits the shape at that key (\ruleName{split-dict}), adding to the residual the
dictionaries that lack the key; once every key of the pattern is in $\beta$ the keys match as a sequence
(\ruleName{pat-dict}), each key $w$ against $\beta(w)$. The empty dictionary pattern matches every
dictionary and leaves the empty residual.
